\section{Introduction}
\label{sec:intro}

Every EEG foundation model begins by deciding what a token carries, and the current generation disagrees: magnitude spectrogram frames~\cite{yang2023biot,pradeepkumar2026tfm}, raw waveform patches~\cite{wang2025cbramod,wang2024eegpt,elouahidi2025reve}, per-bin spectral amplitude and phase~\cite{jiang2024labram}, or discrete codes learned by reconstruction~\cite{ma2025codebrain,sukhbaatar2026bandvq}. What the generation agrees on is the question it does not ask. The labels that matter clinically, seizures, periodic discharges, spike-and-wave, artifacts, are defined by neurophysiologists in a specific vocabulary~\cite{hirsch2021acns}, and no tokenizer paper asks whether its tokens expose that vocabulary. The empirical record suggests the question matters. Two 2026 benchmarks report that small models trained from scratch tie or beat 25M to 69M parameter pretrained models on clinical corpora~\cite{liu2026compass,kontras2026neuroatlas}, and probing studies find that reconstruction-pretrained embeddings under-represent exactly the oscillatory structure clinical events are made of~\cite{kommineni2026aperiodic,tang2026what}. None of these works explains which content a clinical label is a function of, or which part of a model makes that content usable.

The literature treats the choice of token content as a design option to be validated by downstream score. We take the opposite route: the label definitions already specify the content. Read literally, they name quantities that are band-specific, come in two kinds (dynamics of a band's amplitude envelope, and morphology of its waveform), and live at the timescale of the event, 0.1 to 1\,s. A tokenizer either exposes those quantities or leaves the encoder to reconstruct them from something else, and an encoder either reads them along the axes on which the labels are defined or has to learn that organisation from data. Which of these happens should decide the score, and it should decide it in ways that can be predicted before training.

Our answer is the \emph{demodulation hypothesis}. Clinical EEG is well described as a sum of band-limited components, each a slowly varying amplitude envelope modulating a carrier. The labels are low-complexity functionals of two views of this decomposition at the event timescale: the per-band log-envelope (how much, when, how periodically) and the per-band waveform (what shape). A tokenizer that computes these views explicitly, with a learned constant-$Q$ filterbank and an analytic-signal demodulation, raises the \emph{usable} information available to a fixed-capacity encoder in the sense of predictive $\mathcal{V}$-information~\cite{xu2020usable}, although it adds no information in Shannon's sense. An encoder whose attention factorizes along time, electrode and band reads the resulting grid with attention tables of size $P^2 + C^2 + B^2$ instead of $(PCB)^2$, and its advantage should grow as labelled data shrink. Any further per-token transform of the content that carries no label information is a nuisance the reader must learn to undo, and whether it can do so cheaply depends on the reader.

The hypothesis makes predictions that are easy to falsify, and we test them in three stages: controlled signals that isolate each mechanism (Section~\ref{sec:toy}); a progressive attribution on twelve clinical corpora drawn from 425 single-field contrasts, each changing exactly one configuration field, computed from 748 archived training runs (Section~\ref{sec:attribution}); and a comparison against reproduced baselines under corrected protocols (Section~\ref{sec:main}). Our contributions are:
\begin{enumerate}
\item The demodulation hypothesis: a testable account of clinical EEG tokenization with a usable-information argument and eight predictions (Section~\ref{sec:hypothesis}).
\item Evidence that the large, replicated effects in the archive are the predicted ones (band separation, factorized attention, the envelope view for discharges and the waveform view for seizures, sub-second resolution), and that sixteen learned linear rows of the raw patch cannot replace sixteen physical bands under the same encoder (0.135 $\kappa$ on TUEV).
\item A direct measurement of the reader-relativity of usable information: an isometric rotation of the token by a measured cross-frequency direction is neutral in the factorized reader (0.003 over eight paired runs), is harmful when its direction depends on the content (0.03 to 0.24), and moves a flattened host reader by up to 0.2 in either direction.
\item A 1.6M-parameter model that leads the strongest reproduced baselines on TUEV, TUSZ and IIIC, and a CHB-MIT evaluation on the natural window grid, which lowers every model relative to the enriched set used before (0.697 against 0.486 for the same predictions).
\end{enumerate}
